% =============================================================================
% 003 / 068
% Status: DONE
% =============================================================================
% Notes:
%
% Source question:
% 您所整理的吳語播音體語法字尾主要適用的地區是？將來能覆蓋到整個吳語區嗎？
% =============================================================================

\chapter[您所整理的吳語播音體語法字尾主要適用的地區是？將來能覆蓋到整個吳語區嗎？]{您所整理的吳語播音體語法字尾主要適用的地區是？將來能覆蓋到整個吳語區嗎？}
\qaanswer{
基本上適用於吳語、徽語多數地區，是涵蓋南北吳大框架下的語法建構。透過建構，重新梳理了漢字圈各語言的語法脈絡，通宵吳語播音體語法的人甚至可以在此基礎上幫助老湘語、淮東語、贛語、膠遼語甚至漢語普通話進行高階語法建構。題外話，漢語普通話需要嚴謹嚴密的語法梳理，否則這門語言對法學、神學、哲學及自然科學研究是有障礙的。
}

